\documentclass[12pt]{article}
\usepackage[utf8]{inputenc}
\usepackage[T1]{fontenc}
\usepackage{amsmath,amssymb}
\usepackage{geometry}
\geometry{margin=2.5cm}

\begin{document}

\noindent\underline{Lemma 4}: Let $(G,\cdot,e)$ be a finite abelian group.

\[
S=\{\text{ord } a \mid a\in G\} \subseteq \mathbb{N}^*.
\]

\bigskip
\noindent Then:

\begin{enumerate}
\item[(1)] if $m\in S$ and $d\mid m$, $d\in\mathbb{N} \Rightarrow d\in S$.
\item[(2)] $m,n\in S$. \ if $(m,n)=1 \Rightarrow mn\in S$.
\item[(3)] $m,n\in S \Rightarrow [m,n]\in S$.
\end{enumerate}

\bigskip
\noindent\textit{[side note, cited as a general fact]}\footnote{the label in the original (``infinit / abelian'') is unclear -- it appears to note that the fact below holds generally for commuting elements, not just in this particular case.}: \ if $a,b\in G$, with ord$(a)$, ord$(b)$, $ab=ba$, and $\big(\text{ord } a,\text{ord } b\big)=1 \Rightarrow \text{ord}(ab)=\text{ord}(a)\cdot\text{ord}(b)$.

\bigskip
\noindent\big($GL_2(\mathbb{Z})$: 2 matrices of finite order whose product

\noindent has order $\infty$.\big)

\bigskip
\noindent\underline{Proof}:

\bigskip
\noindent (2) \ $\exists\, a,b\in G$ \ ord $a=m$

\noindent ord $b=n$

\[
(ab)^{mn} = (a^m)^n(b^n)^m = e.
\]

\noindent suppose $(ab)^h=e \overset{?}{\Longrightarrow} mn\mid h$.

\[
a^{hn}=b^{-hn}=e
\]
\[
m\mid hn
\]
\[
(m,n)=1
\]
\[
\Rightarrow m\mid h
\]

\noindent similarly $n\mid h$

\[
(m,n)=1 \Rightarrow mn\mid h
\]

\bigskip
\noindent (3)

\noindent $o(a)=m$

\noindent $o(b)=n$

\end{document}
